% Figure 5 of Evans, Hamkins and Perlmutter, Integers 17 (2017), paper G4; typeset from the authors' arXiv source 1510.08155.
\catcode`\.=\active\def.{{1}}
\chessboard[maxfield=m15,
            boardfontsize=8pt,
            label=false,
            labelbottom=true,
            labelleft=true,
            labelfontsize=5pt,
            labelleftwidth=2ex,
            hlabelformat={\hfill\arabic{ranklabel}\,\,},
            showmover=false,
            border=false,
            margin=false,
            % bullet diagonal:
            coloremphstyle=\color{blue!50!black},
            emphfields={a1,b2,c3,d4,e5,f6},
            % bullet exit line:
            pgfstyle=straightmove,
            linewidth=10000sp,
            shortenstart=-2pt,
            shortenend=-20pt,
            color=blue!50!black,
            backmove={g7-a1},
            % bishop bullets:
            coloremphstyle=\color{green!25!black},
            emphfields={h8,i9,j10},
            % bishop safety:
            coloremphstyle=\color{blue!50!black},
            emphfield=g7,
            % defending pawns:
            coloremphstyle=\color{red!50!black},
            emphfields={j2,k3,l4},
            % black threat pawns:
            coloremphstyle=\color{blue!55!red!75!black},
            emphfields={h9,i10,j11},
            % attacking pawns movement:
            pgfstyle={[clockwise=false,style=knight]curvemove},
            linewidth=20000sp,
            shortenstart=0pt,
            shortenend=3pt,
            color=blue!55!red!75!black,
            backmove={h9-i6,i10-j7,j11-k8},
            % defending pawns movement:
            pgfstyle={[clockwise=false,style=knight]curvemove},
            linewidth=20000sp,
            shortenstart=2pt,
            shortenend=5pt,
            color=red!50!black,
            backmove={j2-i6,k3-j7,l4-k8},
            % The position itself:
            setfen=%
.............../%
.........p..../%
........pPp..../%
.......pPbPp/%
......pPbpbP/%
.....pPbpb.p../%
.....Pbpb.pP../%
.....p.b.pP.../%
p...pPb.pP../%
Pp.pP..pP....../%
.PpP..pP...../%
..P..pP....P/%
....pP....P/%
....Pp...P/%
.....Pp../%
......Pp..../%
]%
